\section{Лекция 12}
\subsection{Свойства несобственного интеграла}
\begin{statement}
    Пусть $\Omega$ --- измеримое по Жордану множество, $f\in \mathcal{R}(\Omega)$ и
    \[I=\idotsint\limits_{\Omega}f(x_1,\dots,x_k)\ dx_1\dots dx_k\]
    тогда существует интеграл в новом смысле и они совпадают.
\end{statement}
\begin{proof}
    Пусть $\{K_n\}$ --- исчерпание $\Omega$. Тогда по свойству \ref{podmnoj} интегрируемости на подмножествах существуют интегралы
    \[\idotsint\limits_{K_n}f(x_1,\dots,x_k)\ dx_1\dots dx_k\]
    а тогда по свойству \ref{addit} аддитивности:
    \begin{multline*}
        \left|\idotsint\limits_{K_n}f(x_1,\dots,x_k)\ dx_1\dots dx_k-\idotsint\limits_{\Omega}f(x_1,\dots,x_k)\ dx_1\dots dx_k\right|=\\
        =\left|\idotsint\limits_{\Omega\setminus K_n}f(x_1,\dots,x_k)\ dx_1\dots dx_k\right|\leq M\cdot \mu(\Omega\setminus K_n)
    \end{multline*}
    %Заметим, что $\Omega\setminus K_{n+1}^o\subset\Omega\setminus K_n$. 
    Предположим, что $\mu(\Omega\setminus K_n)\not\to0$, то есть 
    \[\exists\ \epsilon>0: \mu(\Omega\setminus K_n)\geq \epsilon\] 
    Поскольку $\Phi$ --- измеримо по Жордану, то $\forall \epsilon>0$ существуют простые замкнутые $Q_{\epsilon}$ такие, что 
    \[Q_{\epsilon}\subset \Omega,\ \mu(\Omega\setminus Q_{\epsilon})<\frac{\epsilon}{2}\] 
    Заметим, что 
    \[\Omega\setminus K_n=(Q_{\epsilon}\setminus K_n)\sqcup((\Omega\setminus K_n)\setminus Q_{\epsilon})\]
    значит
    \[\mu(Q_{\epsilon}\setminus K_n)=\mu(\Omega\setminus K_n)-\mu((\Omega\setminus K_n)\setminus Q_{\epsilon})> \epsilon-\frac{\epsilon}{2}=\frac{\epsilon}{2}>0\]
    поскольку
    \[\mu((\Omega\setminus K_n)\setminus Q_{\epsilon})\leq \mu(\Omega\setminus Q_{\epsilon})<\frac{\epsilon}{2}\]
    Определим последовательность 
    \[B_n=Q_{\epsilon}\setminus K_{n+1}^o\supset Q_{\epsilon}\setminus K_n\] 
    причем
    \[\mu(Q_{\epsilon}\setminus K_{n+1}^o)\geq\mu(Q_{\epsilon}\setminus K_n)\]
    Получили последовательность $B_n$ непустых замкнутых множеств, таких, что $B_{n+1}\subset B_n$. Значит у $B_n$ есть общая точка $x\in \bigcap\limits B_n$. Тогда $x\in Q_{\epsilon}\subset \Omega$ и $x\not\in K_{n+1}^o$, однако $\bigcup\limits K_n^o=\Omega$ а значит, $x$ должна принадлежать какому-то из $K_n^o$, противоречие. Отсюда $\mu(\Omega\setminus K_n)\to 0$.
\end{proof}
%\noindent Мы доказали, что $\mathcal{R}(\Omega)\subset \widetilde{\mathcal{R}}(\Omega)$.
\begin{statement}
    Пусть $f\in \widetilde{\mathcal{R}}(\Omega)$ и $g$ --- функция, отличающаяся от $f$ только на множестве меры ноль, то есть $\mu(E)=0$, где 
    \[E=\{x\in \Omega:\ f(x)\neq g(x)\}\] 
    причем $g$ --- ограничена на $E$. Тогда $g\in \widetilde{\mathcal{R}}(\Omega)$, причем
    \[\idotsint\limits_{\Omega}g\ dx_1 \dots dx_k=\idotsint\limits_{\Omega}f\ dx_1 \dots dx_k\]
\end{statement}
\begin{proof}
    Пусть $\{K_n\}$ --- исчерпание $\Omega$. По свойству \ref{podmnoj} интегрируемости на подмножестве $\forall n\in \N: f\in \mathcal{R}(K_n)$. Заметим, что $g\neq f$ на множестве $E\cap K_n$ и $\mu(E\cap K_n)=0$ и $g\in \mathcal{B}(K_n)$. Тогда по свойству \ref{diffmery0}:
    \[\idotsint\limits_{K_n}g\ dx_1 \dots dx_k=\idotsint\limits_{K_n}f\ dx_1\dots dx_k\]
    Устремим $n\to \infty$ и получим утверждение.
    %\[\idotsint\limits_{\Omega}g\ dx_1 \dots dx_k=\idotsint\limits_{\Omega}f\ dx_1\dots dx_k\]
\end{proof}
\begin{statement}
    Пусть $f, g\in \widetilde{\mathcal{R}}(\Omega)$. Тогда $\forall c_1, c_2\in \R$:
    \[c_1\cdot f+c_2\cdot g\in \widetilde{\mathcal{R}}(\Omega)\]
    причем
    \begin{multline*}
        \idotsint\limits_{\Omega}(c_1\cdot f+c_2\cdot g)\ dx_1\dots dx_k=\\
        =c_1\cdot \idotsint\limits_{\Omega}f\ dx_1\dots dx_k+c_2\cdot \idotsint\limits_{\Omega}g\ dx_1\dots dx_k
    \end{multline*}
\end{statement}
\begin{proof}
    Пусть $\{K_n\}$ --- исчерпание $\Omega$. По свойствам \ref{podmnoj} и \ref{lin}\\
    $\forall n\in \N: f\in \mathcal{R}(K_n)$ и
    \begin{multline*}
        \idotsint\limits_{K_n}(c_1\cdot f+c_2\cdot g)\ dx_1\dots dx_k=\\
        =c_1\cdot \idotsint\limits_{K_n}f\ dx_1\dots dx_k+c_2\cdot \idotsint\limits_{K_n}g\ dx_1\dots dx_k
    \end{multline*}
    Устремим $n\to \infty$ и получим утверждение.
\end{proof}
\begin{statement}% [Монотонность по функциям] $\\$
    Пусть $f,g\in \widetilde{\mathcal{R}}(\Omega)$ и $f\geq g$ на $\Omega$. Тогда 
    \[\idotsint\limits_{\Omega}f\ dx_1\dots dx_k\geq \idotsint\limits_{\Omega}g\ dx_1 \dots dx_k\]
\end{statement}
\begin{proof}
    Пусть $\{K_n\}$ --- исчерпание $\Omega$. Тогда по свойству \ref{monpofunc}
    \[\idotsint\limits_{K_n}f\ dx_1\dots dx_k\geq \idotsint\limits_{K_n}g\ dx_1\dots dx_k\]
    Устремим $n\to \infty$ и получим утверждение.
\end{proof}
\begin{statement}\label{l12.5}
    Пусть $\Omega_1\subset \Omega_2,\ f\in \widetilde{\mathcal{R}}(\Omega_1)\cap \widetilde{\mathcal{R}}(\Omega_2),\ f\geq 0$ на $\Omega_2$. Тогда
    \[\idotsint\limits_{\Omega_2}f\ dx_1 \dots dx_k\geq \idotsint\limits_{\Omega_1}f\ dx_1 \dots dx_k\]
\end{statement}
\begin{proof}
    Возьмем $\{K_n\}$ --- исчерпание $\Omega_1,\ \{\Phi_j\}$ --- исчерпание $\Omega_2$. Покажем, что $\{K_n\cap \Phi_n\}$ --- также исчерпание $\Omega_1$. Для начала заметим, что $K_n\cap \Phi_n$ --- измеримы по Жордану. Заметим, что
    \[K_n\cap \Phi_n\subset K^o_{n+1}\cap \Phi^o_{n+1}\subset (K_{n+1}\cap \Phi_{n+1})^o\]
    Остается проверить, что 
    \[\bigcup\limits_{n=1}^{\infty}(K_n\cap \Phi_n)=\Omega_1\]
    Для этого покажем, что $\forall m\in \N$:
    \[K_m\subset \bigcup\limits_{n=1}^{\infty}(K_n\cap \Phi_n)\]
    При этом, поскольку $\Omega_1\subset \Omega_2$, то
    \[K_m\subset \Omega_2=\bigcup\limits_{n=1}^{\infty}\Phi_n\]
    \[K_m=K_m\cap\left(\bigcup\limits_{n=1}^{\infty}\Phi_n\right)=\bigcup\limits_{n=1}^{\infty}(K_m\cap \Phi_n)\overset{(1)}{\subset} \bigcup\limits_{n=m}^{\infty}(K_n\cap \Phi_n)\overset{(2)}{=}\bigcup\limits_{n=1}^{\infty}(K_n\cap \Phi_n)\]
    $(1)$: Поскольку при $n\geq m:\ K_m\subset K_n$, а значит $K_m\cap \Phi_n\subset K_n\cap \Phi_n$.\\
    $(2)$: Поскольку $K_n\cap\Phi_n\subset K_{n+1}\cap \Phi_{n+1}$.\\
    Итак, $\{K_n\cap\Phi_n\}$ --- исчерпание $\Omega_1$. Поскольку $K_n\cap \Phi_n\subset \Phi_n$, значит, по свойству \ref{monpopodmnoj} имеем
    \[\idotsint\limits_{\Phi_n}f\ dx_1\dots dx_k\geq \idotsint\limits_{K_n\cap\Phi_n}f\ dx_1 \dots dx_k\]
    Устремим $n\to \infty$ и получим утверждение.
    %\[\idotsint\limits_{\Omega_2}f\ dx_1\dots dx_k\geq \idotsint\limits_{\Omega_1}f\ dx_1 \dots dx_k\]
\end{proof}
% \begin{example}
%     тут будет пример
% \end{example}
% \begin{example}
%     \[\int\limits_{(a,b)}f(x)\ dx\]
%     Если 
%     \[\exists\ \lim\limits_{b'\to b-0}\int\limits_{c}^{b'}f(x)\ dx,\ \exists\ \lim\limits_{a'\to a+0}\int\limits_{a'}^{c}f(x)\ dx\]
%     и они в сумме дают 
%     \[\int\limits_{(a,b)}f(x)\ dx\]
%     Рассмотрим $f(x)$ как на картинке.
%     Известно \[\sum\limits_{n=1}^{\infty}\frac{(-1)^{n+1}}{n}=\ln{2}\],
%     Значит в старом смысле
%     \[\int\limits_{0}^{+\infty}f(x)\ dx=\ln{2}\]
%     но в новом смысле возьмем исчерпание
%     от 0 до 1 и такие же сверху $n^2$ штук, а тех что от 1 до 2 и такие же снизу $n$ штук и там не будет существовать интеграл в новом смысле.
% \end{example}
\subsection{Несобственный кратный интеграл от неотрицательной функции}
\begin{lemma} \label{l12.1}
    Пусть $K\subset \Omega$ --- компакт, $\{K_n\}$ --- какое-либо исчерпание $\Omega$. Тогда $\exists\ n\in \N: K\subset K_n^o$.
\end{lemma}
\begin{proof}
    Из определения знаем, что $K_n\subset K_{n+1}^o$ и
    \[\Omega=\bigcup\limits_{n=1}^{\infty}K_n\subset \bigcup\limits_{n=1}^{\infty}K_{n+1}^o\subset \Omega\]
    значит 
    \[\Omega=\bigcup\limits_{n=1}^{\infty}K_{n+1}^o\]
    Итак, $\{K_{n+1}^o\}$ --- открытое покрытие $K$. Выделим из него конечное подпокрытие $K_{n_1}^o,\dots,K_{n_p}^o$ и получим
    \[K\subset \bigcup\limits_{j=1}^p K_{n_j}^o=K_{n_p}^o\]
\end{proof}
\begin{theorem} [Об одном исчерпании] \label{1ischerp} $\\$
    Пусть $f\geq 0$ на $\Omega$. Следующие условия эквивалентны:
    \begin{enumerate}
        \item $f\in \widetilde{\mathcal{R}}(\Omega)$.
        \item Для любого исчерпания $\{K_n\}$ области $\Omega$, последовательность 
        \[\left\{\idotsint\limits_{K_n}f\ dx_1 \dots dx_k\right\}\]
        ограничена.
        \item Существует такое исчерпание $\{K_n\}$ области $\Omega$, что последовательность 
        \[\left\{\idotsint\limits_{K_n}f\ dx_1 \dots dx_k\right\}\]
        ограничена.
    \end{enumerate}  
\end{theorem}
\begin{proof}\tab
    \begin{enumerate}
        \item $(1)\Rightarrow(2):\ f\in \widetilde{\mathcal{R}}(\Omega)$, значит, существует
        \[I=\lim\limits_{n\to\infty}\idotsint\limits_{K_n}f\ dx_1 \dots dx_k\]
        поэтому для любого исчерпания $\{K_n\}$ искомая последовательность сходится, а значит ограничена.
        \item $(2)\Rightarrow(3)$: Очевидно
        \item $(3)\Rightarrow(1)$: Пусть $\{K'_n\}$ --- исчерпание $\Omega$. Дано, что последовательность
        \[\left\{\idotsint\limits_{K_n}f\ dx_1 \dots dx_k\right\}\]
        ограничена. Заметим, что последовательность монотонно возрастает. Тогда, по теореме Вейерштрасса о монотонной последовательности, существует предел
        \[I=\lim\limits_{n\to\infty}\idotsint\limits_{K_n}f\ dx_1\dots dx_k\]
        По лемме \ref{l12.1}: $\forall m\in \N\ \exists\ n\in \N: K_m'\subset K_n^o$, значит по свойству \ref{podmnoj}: $f\in \mathcal{R}(K_n)\Rightarrow f\in \mathcal{R}(K_m')$. А из свойства \ref{monpopodmnoj} известно, что
        \[\idotsint\limits_{K'_m}f\ dx_1\dots dx_k\leq \idotsint\limits_{K_n}f\ dx_1\dots dx_k\]
        отсюда
        \[\exists\ \lim\limits_{m\to\infty}\idotsint\limits_{K_m'}f\ dx_1 \dots dx_k=J\leq I\]
        Аналогично можно показать, что $I\leq J$.
    \end{enumerate}
\end{proof}
\begin{definition}
Для любого измеримого компакта $K\subset \Omega$
\[\mathcal{R}_{\text{loc}}(\Omega)=\bigcap\limits_{K\subset \Omega} \mathcal{R}(K)\]
То есть это множество всех таких функций $f$, что для любого измеримого компакта $K$ существует
\[\idotsint\limits_{K}f\ dx_1 \dots dx_k\]
Функции $f\in \mathcal{R}_{\text{loc}}(\Omega)$ называются локально интегрируемыми на $\Omega$.
\end{definition}
% пример локально инт но не инт
\begin{example}
    Рассмотрим $f(x)=x$. Заметим, что интеграл
    \[\int\limits_{\R}x\ dx\]
    расходится. Однако для каждого конечного компакта $K$ существует интеграл
    \[\int\limits_{K} x\ dx\]
    Таким образом, $f\in \mathcal{R}_{\text{loc}}(\R)$, но $f\not\in \mathcal{R}(\R)$.
\end{example}
\begin{comm}
    На самом деле для несобсвенных интегралов также верна теорема о интегрируемости на подмножествах (в этом курсе мы ее не рассматриваем). Из этой теоремы следует что $\widetilde{\mathcal{R}}(\Omega)\subset \mathcal{R}_{\text{loc}}(\Omega)$.
\end{comm}
\begin{theorem} [Первый признак сравнения] $\\$
    Пусть $f,g\in \mathcal{R}_{\text{loc}}(\Omega)$ и $0\leq f\leq g$ на $\Omega$. Тогда $g\in \widetilde{\mathcal{R}}(\Omega) \Rightarrow f\in \widetilde{\mathcal{R}}(\Omega)$ и
    \[\idotsint\limits_{\Omega}f\ dx_1 \dots dx_k \leq \idotsint\limits_{\Omega} g\ dx_1 \dots dx_k\]
\end{theorem}
\begin{proof}
    Пусть $\{K_n\}$ --- исчерпание $\Omega$. Тогда 
    \[\idotsint\limits_{K_n}f\ dx_1\dots dx_k\leq \idotsint\limits_{K_n}g\ dx_1\dots dx_k \to I\]
    значит последовательность
    \[\left\{\idotsint\limits_{K_n}f\ dx_1\dots dx_k\right\}\]
    ограничена. Тогда по теореме \ref{1ischerp} об одном исчерпании: $f\in \widetilde{\mathcal{R}}(\Omega)$. Остается устремить $n\to \infty$ и получить
    \[\idotsint\limits_{\Omega}f\ dx_1 \dots dx_k\leq \idotsint\limits_{\Omega}g\ dx_1\dots dx_k\]
\end{proof}
\begin{example} Признак работает только для неотрицательных функций.\\
    Пусть $\Omega=(0,1)$ 
    \[g\equiv 0,\ \ f=-\frac{1}{x}\]
    Заметим ,что $\forall x\in (0,1): f<g$. Существует интеграл 
    \[\int\limits_{(0,1)}0\ dx=0\]
    но не существует
    \[\int\limits_{(0,1)}-\frac{1}{x}\ dx=-\infty\]
\end{example}
\begin{theorem} [Второй признак сравнения] $\\$
    Пусть $f,g\geq 0$ и $\exists\ c_1, c_2>0:\ c_1\cdot f\leq g\leq c_2\cdot f$ на $\Omega$. Тогда интегралы
    \[\idotsint\limits_{\Omega}f\ dx_1\dots dx_k, \ \ \idotsint\limits_{\Omega}g\ dx_1 \dots dx_k\]
    существуют или не существуют одновременно.
\end{theorem}
\begin{proof} По первому признаку сравнения:
    \[g\in \widetilde{\mathcal{R}}(\Omega) \Rightarrow c_1\cdot f\in \widetilde{\mathcal{R}}(\Omega)\Rightarrow f\in \widetilde{\mathcal{R}}(\Omega)\]
    \[f\in \widetilde{\mathcal{R}}(\Omega) \Rightarrow c_2\cdot f\in \widetilde{\mathcal{R}}(\Omega) \Rightarrow g\in \widetilde{\mathcal{R}}(\Omega)\]
\end{proof}
\begin{example}
    Вычислиим интеграл Эйлера Пуассона с помощью кратного интеграла:
    \[I=\int\limits_{-\infty}^{+\infty}e^{-x^2}\ dx\]
    Сначала покажем, что такой интеграл сходится. Заметим, что $f(x)=e^{-x^2}$ --- четная функция. Значит
    \[\int\limits_{-\infty}^{+\infty}e^{-x^2}\ dx=2\cdot\int\limits_{0}^{+\infty}e^{-x^2}\ dx\]
    Заметим, что $e^{-x^2}$ непрерывна на отрезке $[0,1]$, значит интегрируема, а также что $e^{-x^2}\leq e^{-x}$ на отрезке $[1,+\infty]$, при этом
    \[\int\limits_{1}^{+\infty}e^{-x}\ dx=\frac{1}{e}\]
    Значит, по признаку Вейерштрасса интеграл сходится. Рассмотрим двойной интеграл и выберем такое исчерпание $\R^2$:
    \[[-n,n]\times[-n,n]\]
    \begin{multline*}
        \iint\limits_{\R^2}e^{-(x^2+y^2)}\ dxdy=\lim\limits_{n\to \infty}\int\limits_{-n}^{n}\left(\ \int\limits_{-n}^{n}e^{-(x^2+y^2)}\ dy\right)dx=\\
        =\left(\ \int\limits_{-\infty}^{+\infty} e^{-y^2}\ dy\right)\cdot \left(\ \int\limits_{-\infty}^{+\infty}e^{-x^2}\ dx\right)=I^2
    \end{multline*}
    Теперь перейдем к полярным координатам $(r,\phi)$ и рассмотрим рассмотрим исчерпание $\R^2$ кругами: 
    \[r\leq n.\ \phi\in [0,2\pi]\]
    При этом переходе $\det{J}=r$. Значит
    \[\iint\limits_{\R^2}e^{-(x^2+y^2)}\ dxdy=\lim\limits_{n\to \infty}\int\limits_{0}^{2\pi}\left(\int\limits_{0}^{n}r\cdot e^{-r^2}\ dr\right)d\phi=\lim\limits_{n\to \infty}2\pi\cdot \int\limits_{0}^{n}r\cdot e^{-r^2}\ dr\]
    Сделаем замену $r^2=u$ при такой замене $dr=\frac{du}{2r}$ и
    \[\lim\limits_{n\to \infty}2\pi\cdot \int\limits_{0}^{n}r\cdot e^{-r^2}\ dr=\lim\limits_{n\to \infty}2\pi\cdot \frac{1}{2}\cdot \int\limits_{0}^{n^2}e^{-u}\ du=\lim\limits_{n\to \infty}\pi\cdot(1- e^{-n^2})=\pi\]
    Итак $I^2=\pi$, значит $I=\sqrt{\pi}$.
\end{example}
\newpage
